% Fragmento público generado desde propietarios íntegros.
% Residencia: 52.4.2.
% Fuente: ../incoming/radion_monografia_20260827/manuscrito/sections/05_geometria_accion.tex:218-290
\paragraph{Correspondence between phase, action, area, and orientation}

\begin{figure}[H]
\centering
\begin{tikzpicture}[scale=1.12]
  \def\aa{4.0}
  \def\bb{2.95}
  \def\cc{2.70}
  \draw[->,RadGray] (-4.7,0)--(4.7,0) node[right]{\(Q\)};
  \draw[->,RadGray] (0,-3.5)--(0,3.5) node[above]{\(P\)};
  \fill[RadTeal!8] (0,0) ellipse[x radius=\aa,y radius=\bb];
  \draw[RadTeal,line width=1.35pt] (0,0) ellipse[x radius=\aa,y radius=\bb];
  \fill[RadCopper] (-\cc,0) circle (2pt) node[below=2mm]{\(F_-\)};
  \fill[RadCopper] (\cc,0) circle (2pt) node[below=2mm]{\(F_+\)};
  \draw[RadGold,line width=1.1pt,-{Latex}] (\aa,0)
     arc[start angle=0,end angle=57.3,x radius=\aa,y radius=\bb];
  \node[RadGold] at (3.7,2.3) {a radion};
  \draw[dashed,RadBlue] (0,0)--(\aa,0) node[midway,below]{\(a_S\)};
  \draw[dashed,RadBlue] (0,0)--(0,\bb) node[midway,left]{\(b_S\)};
  \node[align=center,RadNavy] at (0,-3.25)
    {total area \(\pi a_Sb_S=2\pi\hret=\hfull\)};
\end{tikzpicture}
\caption{The action ellipse. A parametric arc of one radian sweeps
\(\hret\); one turn encloses \(\hfull\). The foci encode the aperture
\(C^*\) in the distinguished chart.}
\label{fig:elipse-accion}
\end{figure}

\paragraph{Relationship between the reduced action, uncertainty, and area}

The relation \(\hbar=h/(2\pi)\) is often presented as a notational
convenience: when describing oscillations, rotations, and Fourier transforms,
the factors \(2\pi\) disappear if \(\hbar\) is used. The action ellipse
reveals a stronger reading. Division by \(2\pi\) is exactly the
passage from the action of one revolution to the action per unit phase:
\[
h=\mathcal A(2\pi),\qquad \hbar=\mathcal A(1).
\]
The reduced constant is not merely \(h\) in more convenient notation; it is the
areal density of action with respect to phase.

Heisenberg, Dirac, and Jordan established the language of conjugate variables,
commutators, and canonical transformations. No literal
\enquote{ontological ellipse} lacking a source is attributed to them here. The HMT--MD contribution is
to identify the positive figure, construct its semiaxes from \((A,C^*)\), and
prove \eqref{eq:area-theorem}. The historical genealogy opens the language;
the present theorem fixes the object.

\begin{narrativebox}
Action is not appended to motion afterward. In the radion ellipse, action is the area generated by motion. A full turn does not transport a prefabricated constant: it encloses it.
\end{narrativebox}

\Needspace{10\baselineskip}
\paragraph{Algebraic controls of the focal invariants}

Decimal similarities are subject to control:
\begin{center}
\begin{tabularx}{.96\textwidth}{l r X}
\toprule
Comparison & Residue & Verdict\\
\midrule
\(d_1\) versus \(0.54\) & \(4.5455\times10^{-3}\) & not equality; \(54\) is not deduced;\\
\(d_1^2\) versus \(8/27\) & \(2.3352\times10^{-4}\) & proximity without map;\\
\(\eta\) frente a \(3/10\) & \(-3.1032\times10^{-4}\) & no igualdad;\\
\(e_f\) versus \(2/3\) & \(4.7852\times10^{-3}\) & \(2/3\) belongs exactly to \(B_0\);\\
\(e^{-\eta}\) versus \(20/27\) & \(3.0740\times10^{-4}\) & no igualdad.\\
\bottomrule
\end{tabularx}
\end{center}
The robust invariants are \eqref{eq:eta}, \eqref{eq:excentricidad}, and
\eqref{eq:invariante-focal-relativo}. The HMT criterion does not consist in favouring
a decimal resemblance, but in requiring an object, an operation, and conservation.
